% !TeX program = xelatex
\documentclass[12pt]{article}

\usepackage[
  a4paper,
  top=2.4cm,
  bottom=2.4cm,
  left=2.6cm,
  right=2.6cm
]{geometry}
\usepackage{xcolor}
\usepackage{fontspec}
\usepackage{xeCJK}
\usepackage{tikz}
\usepackage{microtype}

\pagestyle{plain}
\setlength{\parindent}{0pt}
\setlength{\parskip}{0.75em}
\linespread{1.18}

\definecolor{awardred}{HTML}{CF1C22}
\definecolor{ink}{HTML}{2B2829}

\IfFontExistsTF{Songti SC}
  {\setCJKmainfont{Songti SC}}
  {\setCJKmainfont{STSong}}
\IfFontExistsTF{PingFang SC}
  {\setCJKsansfont{PingFang SC}}
  {\setCJKsansfont{Heiti SC}}

\newfontfamily\TitleFont[
  Path=fonts/,
  UprightFont=RobotoSlab-Black.ttf
]{RobotoSlab}
\newfontfamily\SansFont[
  Path=fonts/,
  UprightFont=Montserrat-Bold.ttf
]{Montserrat}

\newcommand{\tracking}[2]{\addfontfeatures{LetterSpace=#1}#2}
\newcommand{\scene}[2]{%
  \clearpage
  {\Large\bfseries #1\quad #2}\par
  \vspace{0.6em}
}
\newcommand{\dialogue}[2]{%
  \textbf{#1：}#2\par
}
\newcommand{\stage}[1]{%
  {\sffamily\color{gray}#1}\par
}

\begin{document}

\begin{titlepage}
\begin{tikzpicture}[remember picture, overlay]
  \fill[white] (current page.south west) rectangle (current page.north east);

  \node[
    anchor=north,
    text=ink,
    align=center
  ] at ([yshift=-2.1cm]current page.north) {%
    {\SansFont\fontsize{11pt}{22pt}\selectfont
    \tracking{90}{FOR \hspace{1.2em}YOUR \hspace{1.2em}CONSIDERATION}}%
  };
  

  \node[
    anchor=center,
    align=center,
    text=awardred
  ] at ([yshift=0.7cm]current page.center) {%
    {\TitleFont\fontsize{62pt}{60pt}\selectfont WHO}\\[8pt]
    {\TitleFont\fontsize{62pt}{60pt}\selectfont IS}\\[8pt]
    {\TitleFont\fontsize{62pt}{60pt}\selectfont TAMUD}%
  };

  \node[
    anchor=north,
    align=center
  ] at ([yshift=-21.7cm]current page.north) {%
    {\SansFont\fontsize{20pt}{26pt}\selectfont\textcolor{awardred}{ORIGINAL SCREENPLAY}}\\[8pt]
    {\SansFont\fontsize{17pt}{22pt}\selectfont\textcolor{ink}{Xanchen Zhang}}%
  };
\end{tikzpicture}
\end{titlepage}

\section*{课程大作业}

一个关于当代大学生对 AI 使用限制的思考的剧本创作。

\textbf{主角：}Ta（助教），student\_1--student\_3。

\textbf{片名备选：}绝命 Token Coder / Who Is TaMud / Ta Make U Do It。

\textbf{摘要课程元素：}围绕课程大作业、AI agent 使用限制、学生公平性、学习能力退化与课程规则执行展开。

\scene{Scene-0}{庆功宴，会议室，内}

\dialogue{student\_1（组长）}{我 review 了一下我们的代码，感觉没什么问题，已经提交上去了。虽然规定不让用 AI agent，但是大家这次作业完成得还是不错，大家真是辛苦了。}

\dialogue{s2}{不让用 Claude Code 太抽象了。美国有些地区的大学已经大面积开放 AI 使用了，你老中还在停步不前。}

\dialogue{s3}{我没用过，所以真有传得那么邪乎吗？听说用 Claude Code 这种 agent 写代码，之后就会一次成瘾？之后再也戒不掉了？}

\dialogue{s2}{其实用国产的 Kimi、Mimo 还行，成瘾性不强，偶尔写作业用一用，助教也看不出来。}

\dialogue{s2}{但是外国货才是真爽。Claude Code 的 API，我天天用，感觉助教也看不出来。就是国内有的 API 店家喜欢掺一点其他模型的 API，纯度不够。想用纯度高的 token 还得是去美国买 Claude Max。最近社区里不是很火的那个 BlueCoder 也很棒，美国的一个计算机系老师，叫什么 Water White，带着他的一个学生 Jessy Pinkman 做出来了一个很棒的 agent，叫 BlueCoder。HumanEval 上 pass@1 直接干到 80\%，不是 pass@10，是 pass@1，简直依法入魂，不要太正。}

\dialogue{student\_1}{你收着点吧。再用多一点，你写一个归并排序估计都费劲了。最近助教严抓，我要再看一下你那部分的代码。}

\stage{student\_1 笑骂。}

\dialogue{s2}{怕啥，绝对查不出来。查出来也无所谓啊，Ta 不是你实验室的学长吗？动用一下你的 social 能力啊组长！}

\stage{伏笔。}

\dialogue{student\_1}{诶，这个 TaMud 是谁？}

\stage{s2 和 s3 凑过去看。给面部特写。}

\dialogue{s2 / s3}{TaMud？}

\stage{切到电脑分镜：助教 Ta 弹出消息。}

\dialogue{Ta（消息）}{带着你的组员，带着电脑，一个一个来我的实验室。}

\stage{众人面色铁青。}

\scene{Scene-1}{拷问 s2}

\dialogue{Ta}{s2 带着电脑先进来吧。}

\stage{众人摇头，暗示什么也别说。s2 惊恐。}

\dialogue{Ta}{来说说吧，这次大作业你负责的那一部分？}

\dialogue{s2}{负责后端的数据处理，还有路径搜索那一块。}

\stage{Ta 看屏幕，一边笑，一边点头。}

\dialogue{Ta}{你的代码看上去质量很高啊？是自己写的吗？注释写这么全？}

\dialogue{s2}{都……都是自己写的，助教。}

\dialogue{Ta}{你路径搜索用的什么算法？用的什么数据结构？}

\dialogue{s2}{图……}

\dialogue{Ta}{图用什么存的？}

\stage{s2 慌张沉默，随后逐渐愤怒。}

\dialogue{Ta}{根据课程大纲第 26 条，个人代码 AI 覆盖率超过 10\% 就会直接挂科。你这直接有 99\% 了吧，挂你 100 次都不够。}

\dialogue{s2}{我是用了，但是别的课程早就可以用了啊！为什么就这门课用不了？接纳 AI 就是大势所驱！就这门课搞什么特色课程，要什么公平公正，大家都用不都是公平的吗！}

\dialogue{Ta}{大家都用，就叫公平？}

\dialogue{s2}{不然呢？现在外面都什么环境了？美国那边有的课都开始教学生怎么用 agent 写项目了。人家讲效率，讲工程化，讲 AI-native。我们这边还在问“你是不是自己写的”。这不就是拿旧时代的尺子，量新时代的学生吗？}

\dialogue{Ta}{所以你觉得，会用 Claude Code，就是新时代学生？}

\dialogue{s2}{至少比手搓强。}

\stage{Ta 沉默。给电脑日期特写：6 月 4 日。}

\scene{Scene-2}{拷问组长 student\_1}

\textbf{主要台词：}

\dialogue{Ta}{两年前那个挂着吊针 debug 的学习委员哪里去了？一年前那个为了一个 lab 熬通宵的国奖选手哪里去了？三个月前那个和我一起给同学们兢兢业业写测试的好学弟，到底去哪里了？}

\dialogue{student\_1}{学长，我太想进步了。}

\dialogue{Ta}{看看吧，看看你的手臂都成什么样子了。给我拉下来看看。这些……这些肌肉是一个计算机系的高材生该有的样子吗？从祖师爷图灵开始，哪一个计算机系的大佬是你这个样子！哪里来这么多时间健身？你比我更清楚！}

\end{document}
